\subsection{Introduction}


In Continuous Variable (CV) Quantum Information Theory, one goes beyond discrete systems (such as qubits) to systems described by \textbf{continuous degrees of freedom}
. Many physical systems are naturally continuous. Examples include: Position and momentum of a particle, electromagnetic field modes, vibrational modes of trapped ions.
We traditionally formulate quantum information theory using qubits. Therefore, concepts in CV quantum information are often less familiar and requires - building a unified formalism, which is the main purpose of this review.
\newline
As we mentioned, many realistic quantum systems are continuous. A mode of the electromagnetic field behaves like a \textbf{quantum harmonic oscillator}. Studying the quantization of electromagnetic fields is important, as optical systems are widely used in communication and sensing, considering the experimental accessibility of photons. The theoretical framework we would describe here would be based on: the quantum harmonic oscillator, Creation and annihilation operators, various quantum gates and phase space representations (e.g., Wigner function). The tools built would allow us to describe quantum states in phase space, analyse Gaussian states, and study transformations via linear operations. Ultimately, the formalism developed here would be used for explaining bosonic quantum channels and state discrimination later.



\subsection{Quantization of a Single Electromagnetic Field Mode}

In this section, we will derive the Hamiltonian corresponding to a single mode of the electromagnetic field confined in a cavity. We will show that each mode behaves as a simple harmonic oscillator.

Consider a single mode of the electromagnetic field in a cavity of volume $V = AL$, where $A$ is the cross-sectional area and $L$ is the length along the $z$-direction.

We assume: The field propagates along the $z$-axis, the electric field is polarized along the $x$-direction, the magnetic field is along the $y$-direction.
\\
The electric field is given by:

\begin{equation}
\mathbf{E}(z,t) = \hat{x} \sqrt{\frac{\omega^2}{4 \epsilon_0 V}} \, q(t) \sin(kz)
\end{equation}

The magnetic field is:
\begin{equation}
\mathbf{B}(z,t) = \hat{y} \sqrt{\frac{1}{4 \epsilon_0 V}} \, \dot{q}(t) \cos(kz)
\end{equation}

where:
\begin{itemize}
\item $q(t)$ is the generalized coordinate (field amplitude)
\item $\dot{q}(t) = \frac{dq}{dt}$
\item $\omega = ck$ is the angular frequency
\item $\epsilon_0$ is the permittivity of free space
\end{itemize}

\subsubsection{Hamiltonian of the Field}


The total energy of the electromagnetic field is:
\begin{equation}
H = \frac{1}{2} \int_V dV \left( \epsilon_0 E^2 + \frac{1}{\mu_0} B^2 \right)
\label{eq:energy}
\end{equation}


Substituting $\mathbf{E}$:

\begin{equation*}
E^2 = \frac{\omega^2}{4 \epsilon_0 V} q^2(t) \sin^2(kz)
\end{equation*}

\begin{equation*}
\epsilon_0 E^2 = \frac{\omega^2}{4 V} q^2(t) \sin^2(kz)
\end{equation*}

Now for $\mathbf{B}$:

\begin{equation*}
B^2 = \frac{1}{4 \epsilon_0 V} \dot{q}^2(t) \cos^2(kz)
\end{equation*}

Using $\frac{1}{\mu_0 \epsilon_0} = c^2$:

\begin{equation*}
\frac{1}{\mu_0} B^2 = \frac{c^2}{4 V} \dot{q}^2(t) \cos^2(kz)
\end{equation*}

Since $\omega = ck$:

\begin{equation*}
\frac{1}{\mu_0} B^2 = \frac{\omega^2}{4 V k^2} \dot{q}^2(t) \cos^2(kz)
\end{equation*}

Thus, we have the Total Hamiltonian in eq.~\eqref{eq:energy} as:

\begin{equation*}
H = \frac{1}{2} \int_V dV \left[
\frac{\omega^2}{4 V} q^2(t) \sin^2(kz)
+
\frac{c^2}{4 V} \dot{q}^2(t) \cos^2(kz)
\right]
\end{equation*}


Using:
\begin{equation*}
\int_0^L \sin^2(kz)\,dz = \frac{L}{2}, \quad
\int_0^L \cos^2(kz)\,dz = \frac{L}{2}
\end{equation*}

we obtain:
\begin{equation*}
\int_V \sin^2(kz)\, dV = \frac{V}{2}, \quad
\int_V \cos^2(kz)\, dV = \frac{V}{2}
\end{equation*}

Substituting:

\begin{equation*}
H = \frac{1}{2} \left[
\frac{\omega^2}{4 V} q^2(t) \cdot \frac{V}{2}
+
\frac{c^2}{4 V} \dot{q}^2(t) \cdot \frac{V}{2}
\right]
\end{equation*}

Simplifying:

\begin{equation*}
H = \frac{1}{2} \left[
\frac{\omega^2}{8} q^2(t)
+
\frac{c^2}{8} \dot{q}^2(t)
\right]
\end{equation*}



we rewrite:
\begin{equation}
H = \frac{1}{16} \left[ \omega^2 q^2(t) + c^2 \dot{q}^2(t) \right]
\label{eq}
\end{equation}

We now absorb the constants into the variables by redefining:
\begin{equation*}
q(t) \rightarrow \frac{\omega}{4\sqrt{\hbar\omega}}\, q,
\qquad
\dot{q}(t) \rightarrow \frac{c}{4\sqrt{\hbar\omega}}\, \dot{q}(t)
\end{equation*}

Substituting in eq.~\eqref{eq}, we obtain:
\begin{equation}
H = \hbar \omega (q^2 + \dot{q}^2) = \hbar \omega (q^2 + p^2)
\label{final_form}
\end{equation}
where $p = \dot{q}$

Thus, a single mode of the electromagnetic field behaves as a simple harmonic oscillator with frequency $\omega$. This forms the basis for quantization.




\subsection{Bosonic Modes from the Harmonic Oscillator}

From the previous section, the Hamiltonian of a single electromagnetic field mode can be written as:
\begin{equation*}
\hat{H} = \hbar \omega (\hat{q}^2 + \hat{p}^2)
\end{equation*}
where $\hat{q}$ and $\hat{p}$ are dimensionless quadrature operators.
\subsubsection{Connection to the Standard Harmonic Oscillator}

We now relate the dimensionless operators $(\hat{q}, \hat{p})$ to the physical position and momentum operators $(\hat{x}, \hat{p}_x)$ of the harmonic oscillator of mass m, with oscillation frequency $\omega$.

Define:
\begin{equation}
\hat{q} = \sqrt{\frac{m\omega}{2\hbar}} \, \hat{x}, 
\qquad
\hat{p} = \frac{1}{\sqrt{2m\hbar\omega}} \, \hat{p}_x
\end{equation}

Substituting into the Hamiltonian in eq.~\eqref{final_form}:
\begin{align*}
\hat{H}
&= \hbar \omega \left(
\frac{m\omega}{2\hbar} \hat{x}^2
+
\frac{1}{2m\hbar\omega} \hat{p}_x^2
\right) \\
&= \frac{1}{2} m\omega^2 \hat{x}^2 + \frac{\hat{p}_x^2}{2m}
\end{align*}

Thus, we recover the standard harmonic oscillator Hamiltonian:
\begin{equation*}
\hat{H} = \frac{\hat{p}_x^2}{2m} + \frac{1}{2} m\omega^2 \hat{x}^2
\end{equation*}

In quantum mechanics, the position and momentum operators satisfy:
\begin{equation*}
[\hat{x}, \hat{p}_x] = i\hbar
\end{equation*}

Using the definitions of $\hat{q}$ and $\hat{p}$, we compute:
\begin{align*}
[\hat{q}, \hat{p}]
&= \left[ 
\sqrt{\frac{m\omega}{2\hbar}} \hat{x},
\frac{1}{\sqrt{2m\hbar\omega}} \hat{p}_x
\right] \\
&= \frac{1}{2\hbar} [\hat{x}, \hat{p}_x] \\
&= \frac{1}{2\hbar} (i\hbar) \\
&= \frac{i}{2}
\end{align*}

Thus, the quadrature operators satisfy:
\begin{equation}
[\hat{q}, \hat{p}] = \frac{i}{2}
\label{com}
\end{equation}

\subsubsection{Creation and Annihilation Operators}

We now define the bosonic ladder operators:
\begin{equation*}
\hat{a} = \hat{q} + i\hat{p}, \qquad
\hat{a}^\dagger = \hat{q} - i\hat{p}
\end{equation*}

They satisfy the commutation relation:

\begin{align*}
[\hat{a}, \hat{a}^\dagger]
&= [\hat{q} + i\hat{p}, \hat{q} - i\hat{p}] \\
&= -i[\hat{q}, \hat{p}] + i[\hat{p}, \hat{q}] \\
&= -i\left(\frac{i}{2}\right) + i\left(-\frac{i}{2}\right) \\
&= 1
\end{align*}

\subsubsection{Hamiltonian in Terms of Ladder Operators}

Using:
\begin{align*}
\hat{a}^\dagger \hat{a}
&= (\hat{q} - i\hat{p})(\hat{q} + i\hat{p}) \\
&= \hat{q}^2 + i\hat{q}\hat{p} - i\hat{p}\hat{q} - i^2 \hat{p}^2 \\
&= \hat{q}^2 + \hat{p}^2 + i(\hat{q}\hat{p} - \hat{p}\hat{q}) \\
&= \hat{q}^2 + \hat{p}^2 + i[\hat{q}, \hat{p}] \\
&= \hat{q}^2 + \hat{p}^2 - \frac{1}{2}
\end{align*}

we obtain:
\begin{align*}
\hat{H}
&= \hbar\omega (\hat{q}^2 + \hat{p}^2) \\
&= \hbar\omega \left( \hat{a}^\dagger \hat{a} + \frac{1}{2} \right)
\end{align*}

Physical Interpretation: The operators $\hat{a}$ and $\hat{a}^\dagger$ describe annihilation and creation of quanta of the electromagnetic field. Each mode behaves as a quantum harmonic oscillator, and its energy eigenstates correspond to photon number states.


\subsubsection{Fock States and Vacuum State}

We define the vacuum state $\ket{0}$ as:
\begin{equation}
\hat{a} \ket{0} = 0
\end{equation}

Higher number states are generated by repeated application of the creation operator:
\begin{equation}
\ket{n} = \frac{(\hat{a}^\dagger)^n}{\sqrt{n!}} \ket{0}
\end{equation}

These states satisfy:
\begin{equation}
\hat{N} \ket{n} = n \ket{n}
\end{equation}

where we have defined the number operator as: $\hat{N} = \hat{a}\hat{a}^\dagger$
The annihilation and creation operators act as:

\begin{equation*}
\hat{a}\ket{n} = \sqrt{n}\ket{n-1}, \quad n \geq 1
\end{equation*}

\begin{equation*}
\hat{a}^\dagger \ket{n} = \sqrt{n+1}\ket{n+1}
\end{equation*}

Thus, the Fock states represent states with a definite number of quanta (photons).

\subsubsection{Multi-Mode Bosonic System}

We now generalize to a system of $n$ independent modes.

Each mode is described by its own annihilation operator:
\begin{equation*}
\hat{a}_k, \quad k = 1,2,\dots,n
\end{equation*}

These operators satisfy the canonical commutation relations:
\begin{equation}
[\hat{a}_k, \hat{a}_l^\dagger] = \delta_{kl}
\end{equation}

The Hilbert space of the system is the tensor product of the individual Fock spaces' hamiltonians:
\begin{equation}
\mathcal{H} = \bigotimes_{k=1}^{n} \mathcal{H}_k
\end{equation}

The number operator for each mode is:
\begin{equation}
\hat{N}_k = \hat{a}_k^\dagger \hat{a}_k
\end{equation}

The total Hamiltonian of the system is:
\begin{equation}
\hat{H} = \sum_{k=1}^{n} \hbar \omega_k \left( \hat{a}_k^\dagger \hat{a}_k + \frac{1}{2} \right)
\end{equation}

\subsubsection{Quadrature Operators for Each Mode}

For each mode, we define the quadrature operators:
\begin{equation*}
\hat{q}_k = \frac{1}{2} (\hat{a}_k + \hat{a}_k^\dagger)
\end{equation*}

\begin{equation*}
\hat{p}_k = \frac{1}{2i} (\hat{a}_k - \hat{a}_k^\dagger)
\end{equation*}

These satisfy the commutation relation:
\begin{equation*}
[\hat{q}_k, \hat{p}_l] = \frac{i}{2} \delta_{kl}
\end{equation*}

\subsubsection{Phase Space Representation}

We now introduce the vector of quadrature operators:
\begin{equation*}
\hat{R} = (\hat{q}_1, \hat{p}_1, \hat{q}_2, \hat{p}_2, \dots, \hat{q}_n, \hat{p}_n)^T
\end{equation*}

The commutation relations can be compactly written as:
\begin{equation}
[\hat{R}_i, \hat{R}_j] = \frac{i}{2} \Omega_{ij}
\label{symp}
\end{equation}

where $\Omega$ is the symplectic form:
\begin{equation}
\Omega =
\bigoplus_{k=1}^{n}
\begin{pmatrix}
0 & 1 \\
-1 & 0
\end{pmatrix}
\end{equation}


The symplectic form $\Omega$ can be written explicitly as a $2n \times 2n$ matrix:
\begin{equation}
\Omega =
\begin{pmatrix}
0 & 1 & 0 & 0 & \cdots & 0 \\
-1 & 0 & 0 & 0 & \cdots & 0 \\
0 & 0 & 0 & 1 & \cdots & 0 \\
0 & 0 & -1 & 0 & \cdots & 0 \\
\vdots & \vdots & \vdots & \vdots & \ddots & \vdots \\
0 & 0 & 0 & 0 & \cdots & 1 \\
0 & 0 & 0 & 0 & \cdots & -1
\end{pmatrix}
\end{equation}

This matrix encodes the canonical commutation relations between all quadrature operators and defines the symplectic structure of the phase space.

\subsection{Mean Vector, Expectation Values and the Covariance Matrix}

We now complete the phase-space description by defining expectation values.

For a quantum system described by a density operator $\hat{\rho}$, the expectation value of any operator $\hat{O}$ is defined as:
\begin{equation}
\langle \hat{O} \rangle = \mathrm{tr}(\hat{\rho} \hat{O})
\end{equation}

Here, the density operator $\hat{\rho}$ describes the full state of the system, which may be either pure or mixed.

The mean vector (first moments) of the quadrature operators is defined as:
\begin{equation}
\bar{R}_k = \langle \hat{R}_k \rangle = \mathrm{tr}(\hat{\rho} \hat{R}_k)
\end{equation}

Thus, $\bar{R}$ gives the average values of all quadrature operators and represents the displacement of the state in phase space.

\subsubsection{Covariance Matrix}

The second moments of the quadrature operators are captured by the covariance matrix $\sigma$, defined as:
\begin{equation}
\sigma_{kl}
=
\frac{1}{2}
\langle \hat{R}_k \hat{R}_l + \hat{R}_l \hat{R}_k \rangle
-
\langle \hat{R}_k \rangle \langle \hat{R}_l \rangle
\label{cov_mat}
\end{equation}

Using the anticommutator $\{A,B\} = AB + BA$, this can be written as:
\begin{equation}
\sigma_{kl}
=
\frac{1}{2}
\langle \{ \hat{R}_k, \hat{R}_l \} \rangle
-
\langle \hat{R}_k \rangle \langle \hat{R}_l \rangle
\end{equation}

The covariance matrix describes: the uncertainties (variances) of quadrature operators and the correlations between different modes.

\subsubsection*{Gaussian States}

A quantum state is said to be Gaussian if its phase-space representation (such as the Wigner function) is Gaussian.

In this case, the full quantum state is completely characterized by: the mean vector $\bar{R}$, and the covariance matrix $\sigma$.

When multiple modes are considered, the state is referred to as a joint Gaussian state, meaning that the joint phase-space distribution of all quadratures is a multivariate Gaussian.

\subsubsection*{Physical Significance}

The mean vector determines the position of the state in phase space, while the covariance matrix determines its spread and correlations.

These quantities will become relevant in the study of Gaussian states, Gaussian transformations, and bosonic quantum channels.

The symplectic structure defined by $\Omega$ is further explored in the following section: 


\subsection{The Symplectic Structure of Phase Space and the Group $\mathrm{Sp}(2n,\mathbb{R})$}

In the previous section, we introduced the vector of quadrature operators:
\begin{equation*}
\hat{R} = (\hat{q}_1, \hat{p}_1, \hat{q}_2, \hat{p}_2, \dots, \hat{q}_n, \hat{p}_n)^T
\end{equation*}

These operators satisfy the canonical commutation relations:
\begin{equation*}
[\hat{R}_i, \hat{R}_j] = \frac{i}{2} \Omega_{ij}
\end{equation*}

where $\Omega$ is the \textbf{symplectic form}, given by:
\begin{equation*}
\Omega =
\begin{pmatrix}
0 & 1 & 0 & 0 & \cdots & 0 \\
-1 & 0 & 0 & 0 & \cdots & 0 \\
0 & 0 & 0 & 1 & \cdots & 0 \\
0 & 0 & -1 & 0 & \cdots & 0 \\
\vdots & \vdots & \vdots & \vdots & \ddots & 1 \\
0 & 0 & 0 & 0 & -1 & 0
\end{pmatrix}
=
\bigoplus_{k=1}^{n}
\begin{pmatrix}
0 & 1 \\
-1 & 0
\end{pmatrix}
\end{equation*}

\vspace{0.5em}

\subsubsection{Definition of Symplectic Transformations}

Let us define the set of all real $2n \times 2n$ matrices:
\begin{equation*}
\mathcal{M}(2n) = \{ S \mid S \text{ is a real } 2n \times 2n \text{ matrix} \}
\end{equation*}

A matrix $S \in \mathcal{M}(2n)$ is called \textbf{symplectic} if it satisfies:
\begin{equation}
S^T \Omega S = \Omega
\label{symplectic_condition}
\end{equation}

This condition ensures that the transformation preserves the symplectic structure of phase space.

\vspace{0.5em}

\subsubsection{The Symplectic Group}

The set of all matrices satisfying Eq.~\eqref{symplectic_condition} forms a group under matrix multiplication, called the \textbf{real symplectic group}:
\begin{equation*}
\mathrm{Sp}(2n, \mathbb{R}) = \{ S \in \mathcal{M}(2n) \mid S^T \Omega S = \Omega \}
\end{equation*}

This group has the following properties:
\begin{itemize}
\item Closure: If $S_1, S_2 \in \mathrm{Sp}(2n,\mathbb{R})$, then $S_1 S_2 \in \mathrm{Sp}(2n,\mathbb{R})$
\item Identity: The identity matrix $I$ belongs to $\mathrm{Sp}(2n,\mathbb{R})$
\item Inverse: If $S \in \mathrm{Sp}(2n,\mathbb{R})$, then $S^{-1} \in \mathrm{Sp}(2n,\mathbb{R})$
\item Associativity: Follows from matrix multiplication
\end{itemize}

\vspace{0.5em}

\subsubsection{Real Canonical Linear Transformations}

Consider a real canonical linear transformation (where quadratures obey the canonical commutation relation in eq.~\eqref{symp} of the quadrature operators:
\begin{equation}
\hat{R}' = S \hat{R}
\end{equation}

We now compute the commutation relations of the transformed operators:
\begin{align*}
[\hat{R}'_i, \hat{R}'_j]
&= \left[ \sum_k S_{ik} \hat{R}_k , \sum_l S_{jl} \hat{R}_l \right] \\
&= \sum_{k,l} S_{ik} S_{jl} [\hat{R}_k, \hat{R}_l] \\
&= \sum_{k,l} S_{ik} S_{jl} \left( \frac{i}{2} \Omega_{kl} \right) \\
&= \frac{i}{2} (S \Omega S^T)_{ij}
\end{align*}

As it is a real canonical linear transformation, for the commutation relations to remain invariant, we require:
\begin{equation*}
S \Omega S^T = \Omega
\end{equation*}

Taking transpose of both sides:
\begin{equation}
S^T \Omega S = \Omega
\end{equation}

which is precisely the symplectic condition.

Thus, for every canonical linear transformation S, we have S $\in$ $\mathrm{Sp}(2n,\mathbb{R})$. Conversely for any S $\in$ $\mathrm{Sp}(2n,\mathbb{R})$, we can have a canonical linear transformation $\hat{R}' = S \hat{R}$. Thus, \textbf{real canonical linear transformations in 2n dimensions} and \textbf{real 2n x 2n symplectic matrices} have a one-to-one correspondence. Symplectic transformations (S $\in$ $\mathrm{Sp}(2n,\mathbb{R})$) are important because: they preserve the \textbf{uncertainty relations} in quantum mechanics, and they will describe (as we will see later) \textbf{Gaussian unitary operations} in continuous-variable quantum systems.

\vspace{0.5em}




\subsection{Quadratic Hamiltonians and Gaussian Transformations}

We study Hamiltonians that are at most quadratic in the quadrature operators $\hat{R} = (\hat{R}_1, \dots, \hat{R}_{2n})^T$. The general form is:
\begin{equation}
\hat{H} = \frac{1}{2} \hat{R}^T G \hat{R} + \mathbf{b}^T \hat{R},
\label{eq:quad_ham}
\end{equation}
where $G$ is a real, symmetric $2n \times 2n$ matrix ($G=G^T$) and $\mathbf{b}$ is a real vector. The symmetry of $G$ ensures Hermiticity of our Hamiltonian. Such quadratic Hamiltonians also generate \textit{Gaussian transformations}, mapping Gaussian states to Gaussian states.

\paragraph{Heisenberg Evolution:}

The time evolution of the quadrature operators is governed by the Heisenberg equation. Using the canonical commutation relations $[\hat{R}_i, \hat{R}_j] = \frac{i}{2} \Omega_{ij}$ (eq.~\eqref{symp}), we derive the equation of motion for $\hat{R}$.

Since the quadrature operators $\hat{R}_m$ have no explicit time dependence, the equation of motion for the $m$-th component is:
\begin{equation}
\frac{d\hat{R_m}}{dt} = \frac{i}{\hbar} [\hat{H}, \hat{R_m}].
\end{equation}

First, consider the commutator of the quadratic term of the Hamiltonian $\hat{H} = \frac{1}{2} \hat{R}^T G \hat{R}$ = $\frac{1}{2} \sum_{k,l} G_{kl} \hat{R}_k \hat{R}_l$  with $\hat{R}_m$:
\begin{align}
\left[ \frac{1}{2} \sum_{k,l} G_{kl} \hat{R}_k \hat{R}_l, \hat{R}_m \right] 
&= \frac{1}{2} \sum_{k,l} G_{kl} \left( \hat{R}_k [\hat{R}_l, \hat{R}_m] + [\hat{R}_k, \hat{R}_m] \hat{R}_l \right) \nonumber \\
&= \frac{i}{4} \sum_{k,l} G_{kl} \left( \hat{R}_k \Omega_{lm} + \Omega_{km} \hat{R}_l \right).
\label{27}
\end{align}
By exploiting the symmetry of $G$ and antisymmetry of $\Omega$ ($\Omega^T = -\Omega$), both terms combine into a single matrix product. Specifically, first term in eq.~\eqref{27} $\sum_{k,l} G_{kl} \Omega_{lm} \hat{R}_k = \sum_{k,l} ((G\Omega)^T)_{mk} \hat{R}_k = -\sum_k (\Omega G)_{mk} \hat{R}_k$  (as $\Omega^T = -\Omega$ and $G^T = G$) and similarly for the second term in eq.~\eqref{27}. Thus:
\begin{equation}
[\hat{H}_{quad}, \hat{R}] = -\frac{i}{2} \Omega G \hat{R}.
\end{equation}
Now, for the linear term in H, $[\mathbf{b}^T \hat{R}, \hat{R}_m] = \sum_k b_k \frac{i}{2} \Omega_{km} = -\sum_k\frac{i}{2} \Omega_{mk} b_k  = -\frac{i}{2} (\Omega \mathbf{b})_m$.

Combining these results and applying the Heisenberg equation $\frac{d\hat{R}}{dt} = \frac{i}{\hbar} [\hat{H}, \hat{R}]$ (where constants are absorbed into the definition of $\Omega$), we obtain the linear differential equation:
\begin{equation}
\frac{d\hat{R}}{dt} = \Omega G \hat{R} + \Omega \mathbf{b}.
\label{eq:eom}
\end{equation}


Equation \ref{eq:eom} is a linear differential equation. Defining $A = \Omega G$ and $\mathbf{c} = \Omega \mathbf{b}$, the formal solution is:
\begin{equation}
\hat{R}(t) = S(t)\hat{R}(0) + \mathbf{d}(t),
\label{R}
\end{equation}
where $S(t) = e^{At} = e^{\Omega G t}$ is the homogeneous solution and $\mathbf{d}(t) = \int_0^t e^{\Omega G (t-s)} \Omega \mathbf{b} \, ds$ is the displacement.

\paragraph{Symplectic Nature of $S(t)$:}
We verify that $S(t)$ is a symplectic transformation, i.e., $S(t)^T \Omega S(t) = \Omega$. Differentiating with respect to time:
\begin{equation}
\frac{d}{dt} \left( S^T \Omega S \right) = S^T \left( A^T \Omega + \Omega A \right) S.
\end{equation}
Substituting $A = \Omega G$ and using $G^T=G, \Omega^T=-\Omega$:
\begin{align}
A^T \Omega + \Omega A &= (\Omega G)^T \Omega + \Omega (\Omega G) \nonumber \\
&= G^T \Omega^T \Omega + \Omega^2 G \nonumber \\
&= G (-\Omega) \Omega + \Omega^2 G \nonumber \\
&= -G \Omega^2 + \Omega^2 G \nonumber \\
&= -G (-I) + (-I) G \quad (\text{since } \Omega^2 = -I) \nonumber \\
&= G - G = 0.
\end{align}
Since the time derivative is zero, the quantity $S(t)^T \Omega S(t)$ is constant in time. Evaluating at $t=0$ where $S(0) = I$:
\begin{equation*}
S(t)^T \Omega S(t) = S(0)^T \Omega S(0) = I^T \Omega I = \Omega.
\end{equation*}
Thus, $S(t)$ satisfies the definition of a symplectic matrix for all $t$.







\vspace{0.5em}

\paragraph{Transformation of Mean and Covariance Matrix:}

The mean vector is defined as:
\begin{equation}
\bar{R} = \langle \hat{R} \rangle
\end{equation}

From eq.~\eqref{R}:
\begin{equation}
\bar{R} \rightarrow S \bar{R} + \mathbf{d}
\end{equation}

The covariance matrix is:
\begin{equation}
\sigma_{ij} = \frac{1}{2} \langle \{ \hat{R}_i - \bar{R}_i, \hat{R}_j - \bar{R}_j \} \rangle
\end{equation}


Now, 

\begin{equation}
    \hat{R}' - \bar{R}' = S(\hat{R} - \bar{R}).
\end{equation}

Substituting into the definition of the covariance matrix,
\begin{align*}
    \sigma'_{ij}
    &= \frac{1}{2} \left\langle \left\{ \hat{R}'_i - \bar{R}'_i, \hat{R}'_j - \bar{R}'_j \right\} \right\rangle \\
    &= \frac{1}{2} \left\langle \left\{ (S(\hat{R}-\bar{R}))_i, (S(\hat{R}-\bar{R}))_j \right\} \right\rangle.
\end{align*}

Using linearity,
\begin{equation}
    (S(\hat{R}-\bar{R}))_i = \sum_k S_{ik}(\hat{R}_k - \bar{R}_k),
\end{equation}

we obtain
\begin{align*}
    \sigma'_{ij}
    &= \sum_{k,l} S_{ik} S_{jl}
    \frac{1}{2} \left\langle \left\{ \hat{R}_k - \bar{R}_k, \hat{R}_l - \bar{R}_l \right\} \right\rangle \\
    &= \sum_{k,l} S_{ik} \sigma_{kl} S_{jl}.
\end{align*}

Hence, in matrix form,
\begin{equation}
\sigma \rightarrow \sigma' = S \sigma S^T.
\end{equation}

Under such transformations:
\begin{align}
\sigma &\rightarrow S \sigma S^T
\end{align}


Now, we see how \textbf{Gaussian states} are transformed: 
\newline

The Wigner function for Gaussian states is: 

\begin{equation}
    W(R) = \frac{\exp\left[-\frac{1}{2}(R-\bar{R})^T \sigma^{-1} (R-\bar{R})\right]}{(2\pi)^N \sqrt{\det \sigma}}
\end{equation}

With $R' = SR + d$, we invert to $R = S^{-1}(R' - d)$. The exponent transforms as:

\begin{align*}
    (R-\bar{R})^T \sigma^{-1} (R-\bar{R}) 
    &= (R' - \bar{R}')^T (S^{-1})^T \sigma^{-1} S^{-1} (R' - \bar{R}') \\
    &= (R' - \bar{R}')^T (S \sigma S^T)^{-1} (R' - \bar{R}')
\end{align*}

where $\bar{R}' = S\bar{R} + d$ and $\sigma' = S \sigma S^T$. Since $\det S = 1$, $\det \sigma' = \det \sigma$. Thus:

\begin{equation}
    W'(R') = \frac{\exp\left[-\frac{1}{2}(R'-\bar{R}')^T (\sigma')^{-1} (R'-\bar{R}')\right]}{(2\pi)^N \sqrt{\det \sigma'}}
\end{equation}

Thus, Gaussian states are mapped to Gaussian states, and such unitary transformations are also called \textbf{Gaussian transformations}.







\vspace{0.5em}






\subsection{Characteristic Functions (Quantum)}

The characteristic function of a quantum state $\rho$ is defined as
\begin{equation}
\chi(\alpha) = \mathrm{Tr}\left[\rho D(\alpha)\right], \quad
D(\alpha) = \exp(\alpha a^\dagger - \alpha^* a)
\end{equation}

(for displacement operator $D(\alpha)$ refer to eq.~\eqref{eq:displacementoperator})
Because $a$ and $a^\dagger$ do not commute, different operator orderings give different characteristic functions \cite{braunstein2005cv}, which correspond to different quasi-probability distributions in phase space::
\begin{equation}
\chi(\alpha) = \mathrm{Tr}[\rho D(\alpha)] \quad \text{(Weyl / symmetric)},
\end{equation}
\begin{equation}
\chi_N(\alpha) = \mathrm{Tr}\left[\rho e^{\alpha a^\dagger} e^{-\alpha^* a}\right] \quad \text{(normal ordering)},
\end{equation}
\begin{equation}
\chi_A(\alpha) = \mathrm{Tr}\left[\rho  e^{-\alpha^* a} e^{\alpha a^\dagger}\right] \quad \text{(anti-normal ordering)}.
\end{equation}


% @article{braunstein2005cv,
%   title={Quantum information with continuous variables},
%   author={Braunstein, Samuel L. and van Loock, Peter},
%   journal={Reviews of Modern Physics},
%   volume={77},
%   number={2},
%   pages={513--577},
%   year={2005},
%   publisher={APS}
% }

These correspond respectively to different phase-space distributions (Wigner, $P$, and $Q$ functions).

For Gaussian states,$
\chi(\xi) = \exp\left(
-\frac{1}{2} \xi^T \sigma \xi
+ i \, \xi^T \Omega \bar{R}
\right)$

Here $\xi \in \mathbb{R}^{2n}$ denotes the real phase-space vector corresponding to the complex displacement $\alpha$. So, the state is fully determined by mean $\bar{R}$ and covariance $\sigma$.

\textbf{Relation to Classical Probability:}

In classical theory, the characteristic function of a random variable $x$ is
\begin{equation}
\chi_{ct}(\lambda) = \mathbb{E}[e^{i\lambda x}]
= \int p(x)e^{i\lambda x} dx.
\end{equation}

The quantum version is a direct analogue: as if we are finding expectation value ($\mathrm{Tr}[\rho D(\alpha)$]) of $D(\alpha)$ (like of $e^{i\lambda x}$ in classical theory).



\subsection{Wigner Function}

The Wigner function provides a phase-space representation of a quantum state \cite{weedbrook2012gaussian}. For a single-mode density operator $\hat{\rho}$, it is defined as
\begin{equation}
W(q,p) = \frac{1}{\pi \hbar} \int_{-\infty}^{\infty} dy \,
\langle q - y | \hat{\rho} | q + y \rangle \, e^{2ipy/\hbar}
\end{equation}


% @article{weedbrook2012gaussian,
%   title={Gaussian quantum information},
%   author={Weedbrook, Christian and Pirandola, Stefano and Garc{\'\i}a-Patr{\'o}n, Ra{\'u}l and Cerf, Nicolas J. and Ralph, Timothy C. and Shapiro, Jeffrey H. and Lloyd, Seth},
%   journal={Reviews of Modern Physics},
%   volume={84},
%   number={2},
%   pages={621--669},
%   year={2012},
%   publisher={APS},
%   eprint={1110.3234},
%   archivePrefix={arXiv}
% }

\paragraph{Characteristic Function Representation:}
The Wigner function provides an equivalent description of an $n$-mode quantum state in phase space to the characteristic function. It can be defined as the Fourier transform of the symmetrically ordered characteristic function $\chi(\boldsymbol{\lambda})$:

\begin{equation}
W(\boldsymbol{\alpha}) = \int \frac{d^{2n}\boldsymbol{\lambda}}{\pi^{2n}} \,
\chi(\boldsymbol{\lambda})\,
\exp\!\left( \boldsymbol{\alpha} \boldsymbol{\lambda}^\dagger - \boldsymbol{\alpha}^\dagger \boldsymbol{\lambda} \right)
\label{eq:wigner_char_ft}
\end{equation}

For a single mode, this reduces to
\begin{equation}
W(\alpha) = \int \frac{d^2\boldsymbol{\lambda}}{\pi^2} \,
\chi(\boldsymbol{\lambda})\,
\exp\!\left( \boldsymbol{\alpha} \boldsymbol{\lambda}^\dagger - \boldsymbol{\alpha}^\dagger \boldsymbol{\lambda} \right)
\end{equation}

Here, $\boldsymbol{\alpha} = (\alpha_1, \dots, \alpha_n)^T$ is the vector of complex phase-space coordinates, where each component is related to the quadrature operators by $\alpha_k = (q_k + i p_k)/\sqrt{2}$. The variable $\boldsymbol{\lambda}$ is the integration variable in the Fourier transform, corresponding to the argument of the displacement operator in the characteristic function $\chi(\boldsymbol{\lambda}) = \text{Tr}[\hat{\rho} \hat{D}(\boldsymbol{\lambda})]$.

\paragraph{Properties:}
\begin{itemize}
    \item Normalization: $\int dq\,dp\, W(q,p) = 1$.
    \item $W(q,p)$ is real, but not necessarily positive.
    \item Marginals:
    \[
    \int dp\, W(q,p) = \langle q | \hat{\rho} | q \rangle, \quad
    \int dq\, W(q,p) = \langle p | \hat{\rho} | p \rangle.
    \]
\end{itemize}

\paragraph{Expectation Values:}
\begin{equation}
\langle \hat{A} \rangle = \int d^{2n}\mathbf{R}\,
W_\rho(\mathbf{R})\, A_W(\mathbf{R}),
\end{equation}
where $A_W(\mathbf{R})$ is the Weyl symbol of $\hat{A}$.







\subsection{P Representation}

The Glauber--Sudarshan $P$-distribution provides a diagonal representation of the density operator in the coherent-state basis:
\begin{equation}
\hat{\rho} = \int d^2\alpha \, P(\alpha)\, |\alpha\rangle \langle \alpha|.
\end{equation}

\paragraph{Properties:}
\begin{itemize}
    \item \textbf{Diagonal form:} Expresses $\hat{\rho}$ as a weighted mixture of coherent states.
    \item \textbf{Classicality:} If $P(\alpha)$ behaves like a true probability distribution (positive and regular), the state is classical.
    \item \textbf{Non-classicality:} For many quantum states (e.g., Fock states), $P(\alpha)$ becomes highly singular or non-positive.
    \item \textbf{Characteristic Function:} Related to the normally ordered characteristic function:
    \begin{equation}
    \chi_N(\lambda) = \int d^2\alpha \, P(\alpha)\, e^{\lambda \alpha^* - \lambda^* \alpha}.
    \end{equation}
\end{itemize}


\subsection{Trace Rules in Phase Space}

Trace identities allow operator expressions to be evaluated as integrals in phase space, which are often simpler to compute.

\paragraph{Glauber Expansion:}
Any operator $\hat{O}$ acting on an $n$-mode Hilbert space can be written as
\begin{equation}
\hat{O} = \int \frac{d^{2n}\lambda}{\pi^n} \,
\chi[\hat{O}](\lambda)\, \hat{D}(-\lambda),
\end{equation}
where $\chi[\hat{O}](\lambda)$ is the characteristic function.

\paragraph{Trace Rule (Characteristic Functions):}
For two operators $\hat{O}_1$ and $\hat{O}_2$:
\begin{equation}
\mathrm{Tr}(\hat{O}_1 \hat{O}_2)
= \int \frac{d^{2n}\lambda}{\pi^n} \,
\chi[\hat{O}_1](\lambda)\, \chi[\hat{O}_2](-\lambda).
\end{equation}

\paragraph{Trace Rule (Wigner Functions):}
Equivalently, in terms of Wigner functions:
\begin{equation}
\mathrm{Tr}(\hat{O}_1 \hat{O}_2)
= (\pi)^n \int d^{2n}\lambda \,
W[\hat{O}_1](\lambda)\, W[\hat{O}_2](\lambda).
\end{equation}
which will be used to derive the fidelity and Quantum Chernoff Bound for Gaussian states, in later sections.





\subsection{Uncertainty Principle}

The uncertainty principle follows from the non-commutativity of operators. For any two observables $A$ and $B$, the Robertson uncertainty relation is given by
\begin{equation}
\Delta A \Delta B \geq \frac{1}{2} \left| \langle [A,B] \rangle \right|
\end{equation}

For continuous variable systems, the canonical quadrature operators $\hat{q}$ and $\hat{p}$ satisfy
\begin{equation}
[\hat{q}, \hat{p}] = i\hbar,
\end{equation}
which leads to the standard uncertainty relation
\begin{equation}
\Delta q  \Delta p \geq \frac{\hbar}{2}
\end{equation}

More generally, for an $n$-mode system described by the covariance matrix $\sigma$, the uncertainty principle can be written in compact form \cite{simon1994quantum} as
\begin{equation}
\sigma + \frac{i\hbar}{2}\Omega \geq 0
\end{equation}
where $\Omega$ is the symplectic form.

% @article{simon1994quantum,
%   title={Quantum-noise matrix for multimode systems: U(n) invariance, squeezing, and normal forms},
%   author={Simon, R. and Mukunda, N. and Dutta, B.},
%   journal={Physical Review A},
%   volume={49},
%   number={3},
%   pages={1567--1583},
%   year={1994},
%   publisher={APS}
% }




